% !PRINTED 559
\weberpage{559}

%<<CONT
\[
\begin{aligned}
S_0+S_1-S_2-S_3&=\phantom{-}\frac{\pi^2}{24\omega}-\frac{\pi^2\omega}{24}+\pi i\log f(\omega),\\[4pt]
S_0-S_1-S_2+S_3&=-\frac{\pi^2}{12\omega}-\frac{\pi^2\omega}{24}+\pi i\log f_1(\omega),\\[4pt]
S_0-S_1+S_2-S_3&=\phantom{-}\frac{\pi^2}{24\omega}+\frac{\pi^2\omega}{12}+\pi i\log f_2(\omega).
\end{aligned}
\]

Daraus durch Addition:
\[
3S_0-S_1-S_2-S_3=\pi i\log\sqrt{2},
\]
und wenn man $S_0$ eliminiert:
\[
\begin{aligned}
\tfrac{4}{3}S_1-\tfrac{2}{3}S_2-\tfrac{2}{3}S_3&=\phantom{-}\frac{\pi^2}{24\omega}-\frac{\pi^2\omega}{24}-\frac{\pi i}{3}\log\sqrt{2}+\pi i\log f(\omega),\\[4pt]
-\tfrac{2}{3}S_1-\tfrac{2}{3}S_2+\tfrac{4}{3}S_3&=-\frac{\pi^2}{12\omega}-\frac{\pi^2\omega}{24}-\frac{\pi i}{3}\log\sqrt{2}+\pi i\log f_1(\omega),\\[4pt]
-\tfrac{2}{3}S_1+\tfrac{4}{3}S_2-\tfrac{2}{3}S_3&=\phantom{-}\frac{\pi^2}{24\omega}+\frac{\pi^2\omega}{12}-\frac{\pi i}{3}\log\sqrt{2}+\pi i\log f_2(\omega).
\end{aligned}
\]

Man kann nun ebenso wie vorhin von den Summen $S$ zu
den $S'$ übergehen, in denen $x$, $y$ relativ prim sind, und findet:

\weq{7}{\begin{aligned}
\log f(\omega)&=\tfrac{1}{3}\log\sqrt{2}+\frac{\pi i}{24}\Bigl(-\omega+\frac{1}{\omega}-4S_1'+2S_2'+2S_3'\Bigr),\\
\log f_1(\omega)&=\tfrac{1}{3}\log\sqrt{2}+\frac{\pi i}{24}\Bigl(-\omega-\frac{2}{\omega}+2S_1'+2S_2'-4S_3'\Bigr),\\
\log f_2(\omega)&=\tfrac{1}{3}\log\sqrt{2}+\frac{\pi i}{24}\Bigl(-\omega-\frac{1}{\omega}+2S_1'-4S_2'+2S_3'\Bigr).
\end{aligned}}

Diese Formeln stimmen der Form nach mit \S\,147, (11)
überein, was insofern auffallend ist, als die Art des Grenz\-
überganges beide Male eine verschiedene ist.

\begin{center}
\textbf{\S\,149.\quad Entwickelungen für die Funktion $\log\eta(\omega)$.}
\end{center}

Betrachten wir die Summe

\weq{1}{
U=\sum^{x,y}\frac{1}{(x+\omega y)^{2s}},
}

erstreckt über alle $x$, $y$ mit Ausnahme der Kombination $0$, $0$.

Indem man die Glieder mit $y=0$ absondert, kann man
dafür setzen:

\weq{2}{
U=2\sum_{1,\infty}^{x}\frac{1}{x^{2s}}+2\sum_{1,\infty}^{y}\sum_{-\infty,\infty}^{x}\frac{1}{(x+\omega y)^{2s}};
}

sodann ist
%CONT>>
% [sic] In(7) druckt die f2-Zeile -omega-1/omega statt des aus den vorherigen Zeilen folgenden2omega+1/omega. Quellenformel unveraendert.
% [sic] Der Druck nennt147,(11); die verglichenen drei Logarithmenformeln stehen in147,(12). Beibehalten.
